% IV. ПРОГРАМНА РЕАЛИЗАЦИЯ
% Обща структура, йерархия на класовете, коментар на изходния текст (по А.3).
% Листингите се вмъкват с \lstinputlisting от src/, за да не се преписват.
\section{Програмна реализация}
\label{sec:implementation}

\subsection{Обща структура}

Реализацията е на C++20~\cite{iso_cpp20} и се построява с CMake. Проектът е разделен на
библиотека, която съдържа цялата логика, и изпълним файл, който е тънък слой над нея за
работа от команден ред. Разделението не е украса: модулните тестове се свързват срещу
библиотеката и не минават през командния ред, така че проверката на поведението е
независима от начина на извикване. Публичните заглавни файлове живеят в \code{include/%
blueprint/}, реализацията в \code{src/}, а тестовете в \code{tests/}. Всеки модул от
Раздел~\ref{sec:design} има точно един заглавен файл и точно един файл с реализация
\tabref{tab:files}.

\begin{table}[H]
\caption{Файлове на реализацията и отговорността на всеки}
\label{tab:files}
\centering
\fontsize{11}{13}\selectfont
\begin{tabular}{@{}L{3.9cm}L{8.1cm}@{}}
\toprule
\textbf{Файл} & \textbf{Отговорност} \\
\midrule
\code{json.\{hpp,cpp\}} & Документен модел, четец и записвач на JSON. Съществува, защото изводът на фронтенда е JSON, а проектът не приема външни зависимости. \\
\code{model.\{hpp,cpp\}} & Вътрешният модел и записът му във външен вид. Единственото място, където се пази структура. \\
\code{reader.\{hpp,cpp\}} & Обхождане на синтактичното дърво, извеждане на връзките, изпълнение на компилатора. \\
\code{emit.\{hpp,cpp\}} & Общият обход и двата излъчвателя. \\
\code{diagram.\{hpp,cpp\}} & Прочитане на диаграма обратно в модел, за двете нотации. \\
\code{schema.\{hpp,cpp\}} & Модел към релационна схема, схема към DDL и обратно, схема към заготовки на класове. \\
\code{diff.\{hpp,cpp\}} & Съпоставяне на два модела и отчитане на разликите. \\
\code{behavior.\{hpp,cpp\}} & Модели и излъчватели за случаите на употреба и за дейностите. \\
\code{annotations.hpp} & Макросите за анотиране, само дефиниции на препроцесора. \\
\code{io.\{hpp,cpp\}} & Четене и запис на файл, за да се отказват еднакво навсякъде. \\
\code{main.cpp} & Команден ред. Не взема решения, само подрежда извиквания. \\
\bottomrule
\end{tabular}
\end{table}

\subsection{Четене на изходния текст}

Четенето става чрез изпълнение на компилатора като външен процес и разбор на изведеното
от него синтактично дърво в JSON~\cite{clang_ast_json}. Съставянето на командата е
отделено в \code{build\_clang\_command}, което я прави проверима без да се изпълнява
процес, а разборът на резултата е отделен в \code{read\_ast\_json}, което е чиста функция
без достъп до файловата система. Разделянето не е формално: то позволява целият четец да
се тества върху записан образец на дърво, без инсталиран компилатор.

Три подробности от обхождането заслужават коментар, защото и трите са места, на които
наивната реализация мълчаливо произвежда грешен модел.

\textbf{Проследяване на текущия файл.} Фронтендът извежда полето с името на файла
\term{само когато то се променя}, а следващите възли го наследяват. Затова текущият файл
се пренася през цялото обхождане в реда на документа. Пропускането на тази подробност е
разликата между модел на проекта и модел на всеки системен заглавен файл, който проектът
включва.

\textbf{Отсяване на подразбиращите се членове.} Фронтендът извежда и членовете, които
самият той е породил: подразбиращите се конструктори, операторите за присвояване,
деструкторът, както и вграденото име на класа в самия него. Всички те носят признака
\code{isImplicit} и се отхвърлят. Без това всеки клас в диаграмата получава по пет
операции, които никой не е писал.

\textbf{Видимост по секции.} Видимостта на един член не се извежда като негово поле, а
следва от предхождащия го спецификатор за достъп. Затова тя се води при обхождането,
започвайки от \code{private} за \code{class} и от \code{public} за \code{struct}.

Разрешаването на имената до връзки се извършва \term{след} обхождането, когато всички
класове вече са известни, защото типът на член може да сочи към клас, деклариран по-късно
в същата единица за превод. Съответствието на име се прави първо по напълно квалифицирано
име, а ако то не съвпадне, по последната съставка. Двусмислено име не се разрешава до
предположение, а остава неразрешено и не поражда връзка.

\subsection{Заместващи заглавни файлове при извличането}

Извличането се изпълнява с \code{-nostdinc++} и с папката \code{stdstub/} като системна
папка за включване. Тя съдържа заглавни файлове само с декларации, заместващи тези от
стандартната библиотека, които заглавните файлове на проекта включват. Причината е
измерена, а не стилистична: една единица за превод, която включва действителните
\code{<memory>}, \code{<vector>} и \code{<string>}, извежда около 280 MB JSON, почти
изцяло вътрешности на стандартната библиотека, които никоя диаграма на класове не би
показала. Със заместващите файлове същата единица извежда под мегабайт. Те никога не се
компилират в инструмента и никога не се свързват, а извличането се изпълнява с
\code{-fsyntax-only}.

\subsection{Излъчватели}

Излъчвателите в двете нотации~\cite{plantuml_guide,mermaid_docs} споделят един обход на
модела и се различават само по низовете, които произвеждат за всеки елемент. Общата част
е изнесена в базов клас с виртуални методи за отделните елементи, а конкретните
излъчватели предефинират само тях. Редът на елементите и решението кое си заслужава да
бъде излъчено живеят в базовия клас, веднъж.

Това е и мястото, където се вижда цената на разликите между нотациите. Mermaid чете
ъгловите скоби като свой собствен ограничител за параметризация и използва тилда, така че
всяко изписване на тип се превежда на излизане. Нито една от двете нотации няма място за
етикетирана стойност или за списък на шаблонните параметри, затова те пътуват в
\term{бележка}. Бележката е механизмът на самия UML за пояснение върху елемент, тоест
това не е скрит канал, а предвиденият. Един признак остава без превод: Mermaid отбелязва
операция като абстрактна или като статична и няма нищо за операция, която е виртуална, но
не е чиста. Загубата е измерена, а не премълчана, и е отчетена в Раздел~\ref{sec:results}.

Четецът на диаграми приема диалекта, който самият Blueprint произвежда. Ограничението е
съзнателно и е заявено като такова: целта на четенето е да затвори цикъла и да даде на
проверяващия модул модел на документираното решение, а не да приеме всяка диаграма, която
човек би могъл да нарисува.

\subsection{Пораждаща посока}

Пораждащата посока е реализирана като две отделни преобразувания над един и същ модел.
Първото произвежда описание на релационна схема, второто заготовки на класове. И трите
стратегии за наследяване са реализирани, като изборът се прави от етикетирана стойност
върху базовия клас и само при липсата ѝ от подадената от командния ред стойност. При
стратегията с една таблица всяка колона, дошла от наследник, се записва като допускаща
празна стойност, защото ред от друг наследник я оставя празна, и се добавя колона
разграничител. При стратегията с таблица на клас първичният ключ на наследника е и външен
ключ към базата. При стратегията с таблица на конкретен клас абстрактният клас не получава
таблица, а наследените колони се копират надолу.

Породеният DDL носи името на класа в коментар пред всяка таблица и името на атрибута в
коментар след всяка колона, чието име се различава от неговото. Това не е украса.
Възстановяването на име на клас от име на таблица би означавало отгатване на английско
единствено число, а отгатването е точно това, което един пълен цикъл не бива да съдържа.

Заготовките се пораждат само за декларации, тоест инструментът не пише тела на операции и
не докосва съществуващ файл, който вече съдържа реализация. Породеният заглавен файл носи
и макросите за анотиране, така че повторното извличане от него дава съпоставим модел, с
което цикълът се затваря. Че породеният файл е действителен C++, се проверява в
демонстрацията чрез \code{clang -fsyntax-only} върху него, а не се твърди.

\subsection{Проверяващ модул}

Проверяващият модул сравнява модел, извлечен от кода, с модел, прочетен от диаграма, и
извежда разликите като списък. Видовете разлика са девет: добавен, премахнат и променен,
поотделно за класификатор, член и връзка. Съответствието на връзките се прави по двата
края и по етикета, така че промяна на вида или на кратността се отчита като промяна, а не
като премахване плюс добавяне. Разликата е съществена: два реда вместо един за една и съща
промяна биха направили отчета по-шумен точно там, където той трябва да е най-полезен.

Посоката на думите е фиксирана: първият подаден модел е кодът, вторият е решението, така
че клас, който е в кода, но не е в диаграмата, е \term{добавен}, а обратното е
\term{премахнат}. Процесът завършва с код 0 при съвпадение и с код 1 при разлика, което е
и условието проверката да може да прекъсне построяване, вместо да чака някой да я
прочете.

\subsection{Модулни тестове}

Модулните тестове са 58 в 7 групи и не използват външна рамка. Изпълнителят е един
заглавен файл с макроси за твърдение и регистър, попълван при статично инициализиране;
групите се регистрират в CTest поотделно с \code{add\_test}, така че CTest съобщава по
един ред на група, а всяка група изброява своите случаи. Тестовете на четеца работят
върху записан образец на синтактично дърво, в който присъстват само полетата, които
четецът действително чете, тоест образецът е едновременно вход на теста и заявление какво
четецът изисква от фронтенда.
